\documentclass[11pt,a4paper]{article}

% --- Encoding / Sprache ---
\usepackage[utf8]{inputenc}
\usepackage[T1]{fontenc}
\usepackage[ngerman]{babel}

% --- Mathe / Physik ---
\usepackage{amsmath,amssymb,amsfonts,amsthm}
\usepackage{mathtools}
\usepackage{physics}

% --- Layout / Links ---
\usepackage{geometry}
\geometry{margin=2.5cm}
\usepackage{hyperref}
\usepackage{microtype}

% --- Figures / Tables (optional) ---
\usepackage{graphicx}
\usepackage{booktabs}

% --- Theorem environments ---
\newtheorem{definition}{Definition}
\newtheorem{remark}{Bemerkung}

% --- Notation ---
\newcommand{\Z}{\mathbb{Z}}
\newcommand{\C}{\mathbb{C}}
\newcommand{\modtwelve}{\ (\mathrm{mod}\ 12)}
\newcommand{\wT}{\omega_{12}}
\newcommand{\EE}{\mathrm{E}}
\newcommand{\AA}{\mathrm{A}}
\newcommand{\BB}{\mathrm{B}}
\newcommand{\CC}{\mathrm{C}}

\title{\textbf{Diskrete Messklassen aus Record-Holonomie:}\\
	\textbf{Rekonstruktion des Kohärenzfaktors $\gamma$ mit $N{=}12$-Phasenraster und Defekt-Schalter $R\in\Z_{12}$}}
\author{(Entwurf)}
\date{\today}

\begin{document}
	\maketitle
	
	\begin{abstract}
		Wir formulieren ein experimentell testbares Mess-Framework, in dem der nichtselektive Kohärenzfaktor
		$\gamma$ eines Qubit-Systems $S$ nicht phänomenologisch als „Dekohärenzrate“ eingeführt wird, sondern
		als \emph{instrumentell rekonstruierbare} Größe $\gamma=\langle r_1|r_0\rangle$ aus der Überlappung zweier
		Record-Zustände. Der Recordraum wird als Ququart mit vier Kanälen $\{\EE,\AA,\BB,\CC\}$ realisiert, deren
		Phasen digital auf ein $N{=}12$-Raster quantisiert werden. Eine globale Holonomie-/Defektklasse
		$R\in\Z_{12}$ wird als Nebenbedingung $K_\EE+K_\AA+K_\BB+K_\CC\equiv R$ implementiert und wirkt als
		„Schalter“, der \emph{diskrete Messklassen} in der komplexen Ebene erzeugt. Wir geben zwei konkrete
		Blueprints (Photonik und supraleitende Schaltungen), ein Minimalprotokoll zur Rekonstruktion von $\gamma$
		aus $\langle\sigma_x\rangle,\langle\sigma_y\rangle$ ohne volle Tomografie, sowie eine Datenpipeline
		(Clustering, Jitter-Modell, Nulltest). Messbar sind u.\,a. Plateau-/Cluster-Strukturen von $\gamma$,
		sprunghafte Wechsel unter $R$-Sweep sowie topologisch selektive Minimierung $|\gamma|\approx 0$ im
		Gegenpaar-Muster $\{u,v,u{+}6,v{+}6\}$.
	\end{abstract}
	
	\tableofcontents
	
	\section{Einführung und Motivation}
	
	Messrückwirkung und Dekohärenz werden in vielen Plattformen durch effektive Parameter beschrieben
	(z.\,B. „Messstärke“, „Dephasingrate“). In dieser Arbeit schlagen wir eine komplementäre
	Operationalisierung vor: Der nichtselektive Kohärenzverlust eines Systems $S$ wird über die
	\emph{Record-Geometrie} bestimmt. Konkret wird der Kohärenzfaktor
	\[
	\rho_{01}\ \mapsto\ \gamma\,\rho_{01}
	\]
	als Überlappung zweier Recordzustände $\gamma=\langle r_1|r_0\rangle$ modelliert und experimentell aus
	nur zwei Observablen \(\langle\sigma_x\rangle,\langle\sigma_y\rangle\) rekonstruiert.
	
	Der zentrale Zusatz gegenüber Standardinstrumenten besteht in (i) einer \emph{digitalen}
	Phasenquantisierung auf \(N{=}12\) und (ii) einem globalen \emph{Holonomie-/Defektparameter} \(R\in\Z_{12}\),
	der die zulässigen Viererkombinationen (Kanalindizes) über eine Nebenbedingung koppelt und so
	\emph{klassenweise} Messrückwirkung erzeugt.
	
	\section{Instrumentmodell und Definition des Kohärenzfaktors}
	
	\subsection{Stinespring-Dilation und nichtselektive Dynamik}
	
	Wir betrachten ein Qubit-System \(S\) mit computational basis \(\{\ket0,\ket1\}\) und einen Recordraum \(R\).
	Das Messinstrument wird als Dilation formuliert:
	\begin{equation}
		\ket0_S \mapsto \ket0_S\otimes \ket{r_0},\qquad
		\ket1_S \mapsto \ket1_S\otimes \ket{r_1},
		\label{eq:dilation}
	\end{equation}
	wobei \(\ket{r_0},\ket{r_1}\in\mathcal{H}_R\) nicht notwendig orthogonal sind.
	Wird der Record nicht ausgelesen (nichtselektiv), ergibt sich für die reduzierte Dichtematrix
	\(\rho_S=\mathrm{Tr}_R(\rho_{SR})\):
	\begin{equation}
		\rho_{01}\ \mapsto\ \langle r_1|r_0\rangle\,\rho_{01}.
	\end{equation}
	
	\begin{definition}[Instrumenteller Kohärenzfaktor]
		Der (nichtselektive) Kohärenzfaktor des Instruments ist
		\begin{equation}
			\boxed{\ \gamma \coloneqq \langle r_1|r_0\rangle\ \in\C\ }.
			\label{eq:gamma_def}
		\end{equation}
	\end{definition}
	
	\subsection{Direkte Rekonstruktion von \(\gamma\) ohne volle Tomografie}
	
	Bereite \(S\) in \(\ket{+}=\tfrac1{\sqrt2}(\ket0+\ket1)\) vor, wende das Instrument an, ignoriere \(R\),
	und messe am System:
	\[
	s_x=\langle\sigma_x\rangle,\qquad s_y=\langle\sigma_y\rangle.
	\]
	Dann gilt (bis auf eine konstante Referenzrotation \(e^{-i\phi_0}\), kalibrierbar):
	\begin{equation}
		\boxed{\ \gamma_{\mathrm{meas}} = (s_x+i s_y)\,e^{-i\phi_0}\ }.
		\label{eq:gamma_meas}
	\end{equation}
	
	\section{Vierkanal-Record, $N{=}12$-Raster und Holonomie-Schalter $R$}
	
	\subsection{Record als Ququart mit vier Kanälen}
	
	Wir wählen \(\mathcal{H}_R=\mathrm{span}\{\ket{\EE},\ket{\AA},\ket{\BB},\ket{\CC}\}\) und definieren
	Recordzustände als kanalgewichtete Superpositionen mit Phasen:
	\begin{equation}
		\ket{r_x}=\sum_{i\in\{\EE,\AA,\BB,\CC\}} \sqrt{w_i}\,e^{i\,\theta^{(x)}_i}\ket{i},
		\qquad w_i\ge 0,\ \sum_i w_i=1.
	\end{equation}
	Dann ist
	\begin{equation}
		\gamma=\sum_{i} w_i\,e^{i(\theta^{(0)}_i-\theta^{(1)}_i)}.
		\label{eq:gamma_general}
	\end{equation}
	
	\subsection{Digitales Phasenraster $N{=}12$}
	
	Wir quantisieren die (relativen) Phasen auf ein 12er Raster:
	\begin{equation}
		\wT \coloneqq e^{2\pi i/12},\qquad
		e^{i(\theta^{(0)}_i-\theta^{(1)}_i)}=\wT^{K_i},\quad K_i\in\Z_{12}.
	\end{equation}
	Im symmetrischen Referenzfall \(w_i=\tfrac14\) folgt:
	\begin{equation}
		\boxed{\ \gamma_{\mathrm{pred}}=\frac14\left(\wT^{K_\EE}+\wT^{K_\AA}+\wT^{K_\BB}+\wT^{K_\CC}\right)\ }.
		\label{eq:gamma_pred_sym}
	\end{equation}
	Damit nimmt \(\gamma\) nur diskret viele Werte an (Phasorsumme aus 12-ten Einheitswurzeln).
	
	\subsection{Holonomie-/Defektklasse $R\in\Z_{12}$ als Nebenbedingung}
	
	Die globale Klasse \(R\) wird als Constraint implementiert:
	\begin{equation}
		\boxed{\ K_\EE+K_\AA+K_\BB+K_\CC \equiv R \modtwelve\ }.
		\label{eq:constraint_R}
	\end{equation}
	Operational lässt man drei Indizes frei und bestimmt den vierten:
	\begin{equation}
		\boxed{\ K_\CC \equiv R - K_\EE - K_\AA - K_\BB \modtwelve\ }.
		\label{eq:KC}
	\end{equation}
	
	\begin{remark}[Interpretation als „Holonomie-Schalter“]
		Obwohl \eqref{eq:constraint_R} programmatisch umgesetzt werden kann, wirkt \(R\) wie eine globale,
		nichtlokale Phasenklasse: Für feste \((K_\EE,K_\AA,K_\BB)\) erzeugt ein \(R\)-Sweep eine diskrete Trajektorie
		\(\gamma(R)\) in der komplexen Ebene (sprunghafte Klassenwechsel).
	\end{remark}
	
	\section{Experimentelle Realisierung: Photonik-Blueprint}
	
	\subsection{Hardware-Minimalaufbau}
	
	\begin{itemize}
		\item \textbf{System \(S\):} Photonen-Qubit (Polarisation oder Pfad).
		\item \textbf{Record \(R\):} vier räumliche Modi (Waveguides) als \(\ket{\EE},\ket{\AA},\ket{\BB},\ket{\CC}\).
		\item \textbf{Gewichte \(w_i\):} Koppler-/MZI-Verhältnisse; Referenz: \(w_i=1/4\).
		\item \textbf{Phasen \(\Phi_i\):} pro Arm programmierbarer Phasenschieber; digitale 12-Stufen-Kalibration.
		\item \textbf{Holonomie \(R\):} Implementierung über \eqref{eq:KC} als „Constraint-Wiring“ (C-Phase aus
		Tripel+R).
	\end{itemize}
	
	\subsection{Messprotokoll (minimal, tomografiefrei)}
	
	\begin{enumerate}
		\item Präpariere \(S\) in \(\ket{+}\).
		\item Wende das Instrument \eqref{eq:dilation} an (Kopplung \(S\leftrightarrow R\)).
		\item Traciere/ignoriere \(R\) (keine Record-Detektion oder Summe über alle Ausgänge).
		\item Messe am System \(s_x=\langle\sigma_x\rangle\) und \(s_y=\langle\sigma_y\rangle\).
		\item Rekonstruiere \(\gamma_{\mathrm{meas}}\) via \eqref{eq:gamma_meas}.
	\end{enumerate}
	
	\subsection{Kalibration}
	
	\begin{itemize}
		\item \textbf{Phase-Calibration:} pro Arm DAC\(\rightarrow\)Winkel \(\Phi=\frac{2\pi}{12}K\).
		\item \textbf{Gewichte:} Intensitätsmessung \(\Rightarrow w_i\).
		\item \textbf{Globalphase:} Referenzsetting \(K_i=0\) liefert \(\phi_0\) zur Korrektur in \eqref{eq:gamma_meas}.
	\end{itemize}
	
	\section{Experimentelle Realisierung: Supraleitende-Blueprint}
	
	\subsection{Grundidee}
	
	\begin{itemize}
		\item \textbf{System \(S\):} Transmon-Qubit.
		\item \textbf{Record \(R\):} zwei zusätzliche Qubits (4-dimensional) oder höhere Transmon-Level als Qudit.
		\item \textbf{Instrument:} gezielt implementierte generalisierte Messung (POVM/Kraus-Operatoren), wahlweise
		im Ensemblemodus (wie Photonik) oder im Trajektorienmodus (Record wird ausgelesen).
	\end{itemize}
	
	\subsection{Holonomie-Schalter}
	
	Als technischer Knopf bietet sich Flusssteuerung \(\Phi\) in Ring-/Interferenz-Geometrien an, die effektiv
	phasenartige Offsets in kontrollierten Couplings erzeugt. Operational genügt die digitale Rasterung der
	entsprechenden Phase auf 12 Werte und die constraint-basierte Festlegung von \(K_\CC\) gemäß \eqref{eq:KC}.
	
	\section{Messkampagnen und Datenpipeline}
	
	\subsection{Kampagnenliste und Vorhersage}
	
	Für jede Kampagne \((K_\EE,K_\AA,K_\BB,R)\) wird \(K_\CC\) durch \eqref{eq:KC} bestimmt. Gemessen wird
	\(\gamma_{\mathrm{meas}}\) über \eqref{eq:gamma_meas}, vorhergesagt \(\gamma_{\mathrm{pred}}\) über
	\eqref{eq:gamma_general} bzw. \eqref{eq:gamma_pred_sym}.
	
	\subsection{Postprocessing: Clustering und Fehler}
	
	\begin{itemize}
		\item Plot: \(\gamma_{\mathrm{meas}}\) in der komplexen Ebene, gruppiert nach \(R\) bzw. nach Tripeln.
		\item Fehlermaß:
		\[
		\Delta \coloneqq \abs{\gamma_{\mathrm{meas}}-\gamma_{\mathrm{pred}}}.
		\]
		\item Clustering: Distanzschwelle \(\varepsilon\) aus Phasenjitter; Report: Clusteranzahl pro \(R\),
		Zentren, Varianzen, Trajektorien \(\gamma(R)\).
	\end{itemize}
	
	\subsection{Jitter-Modell (linearisiert)}
	
	Phasenrauschen \(\Phi_i\to \Phi_i+\delta\Phi_i\) führt zu
	\[
	\delta\gamma \approx i\sum_i w_i\,\wT^{K_i}\,\delta\Phi_i,
	\]
	woraus sich Clusterbreiten und erwartete Punktwolkenform (annähernd elliptisch) abschätzen lassen.
	
	\section{Fünf messbare Signaturen}
	
	\begin{enumerate}
		\item \textbf{\(\gamma\)-Klassierung (Plateaus/Cluster):} bei 12er Raster clustern \((\langle\sigma_x\rangle,\langle\sigma_y\rangle)\)
		auf diskreten Punkten.
		\item \textbf{\(R\)-Sprünge:} für fixes \((K_\EE,K_\AA,K_\BB)\) springt \(\gamma(R)\) beim Umschalten \(R\mapsto R+\Delta R\).
		\item \textbf{Topologisch selektiver Nulltest:} gewisse Muster erlauben \(|\gamma|\approx 0\), andere nicht.
		\item \textbf{Rückgewinnbarkeit (Quantenradierer, optional):} partielle Recordmessung in alternativer Basis
		stellt Kohärenz abhängig von \(|\gamma|\) wieder her.
		\item \textbf{Messklassen statt Messrate (MIPT-Anschluss, optional):} in Plattformen mit variabler Messung
		kann die Klasse (diskretes \(\gamma\)-Spektrum + \(R\)-Schalter) als Kontrollparameter dienen.
	\end{enumerate}
	
	\section{Nulltest: Gegenpaar-Muster $\{u,v,u{+}6,v{+}6\}$}
	
	Im symmetrischen Fall \(w_i=\tfrac14\) gilt \(\gamma=0\) genau dann, wenn die vier Phasoren paarweise
	gegenläufig sind:
	\begin{equation}
		\boxed{\ \gamma=0\ \Longleftrightarrow\ \{K_\EE,K_\AA,K_\BB,K_\CC\}=\{u,v,u+6,v+6\}\modtwelve\ }.
	\end{equation}
	Experimentell prüfbar durch gezielte Settings mit kleinem Fringe-Kontrast und dessen „Aufhebung“ durch \(R\)-Schaltung.
	
	\section{Diskussion und Einordnung}
	
	Das vorgeschlagene Framework verschiebt den Fokus von kontinuierlichen Dekohärenzparametern hin zu
	\emph{instrumentell kodierten Klassen} der Messrückwirkung. Die Kombination aus (i) digitaler
	Phasenquantisierung, (ii) Vierkanal-Record und (iii) globalem Holonomie-Schalter \(R\) liefert ein
	konkretes, laborfähiges Testbett für „Dekohärenzspektren“ (diskret) statt reiner „Dekohärenzraten“ (kontinuierlich).
	Die primäre Aussage ist operational: \(\gamma\) ist messbar und schaltbar, und seine Punktstruktur in der
	komplexen Ebene ist eine direkte Signatur der implementierten Recordgeometrie.
	
	\section*{Danksagung / Hinweis}
	Dies ist ein konzeptioneller Entwurf mit klaren Messsignaturen; eine vollständige physikalische Identifikation
	des zugrunde liegenden Instruments (Kraus-Operatoren, Hardwaredetails, Fehlerbudget) kann direkt auf
	dieser Vorlage aufbauen.
	
\end{document}
